\chapter{Vectorised Backtests}\label{ch:rs:vectorised-backtests}

A one-line \omterm{def:rs:vectorised-backtests:backtest}{backtest} of the one-day reversal on the synthetic market (buy yesterday's losers, sell its winners, every day) shows a Sharpe
ratio of 7.2 over ten years. Written one line differently, multiplying the book built from today's close by today's return, it shows
$-65$. Neither is a result. The first trades at the closing price it used to compute the signal, ignores trading costs on a book that
turns over three quarters of itself every day, and would lose money in any of the ways this chapter lists; the second is a look-ahead so
blatant that it inverts the signal. The \omterm{def:rs:vectorised-backtests:backtest}{vectorised backtest}, a few array operations on a panel of weights and returns, is the fastest
research tool the firm has, and most of the time the right one; this chapter builds it (\texttt{firm.vecbt}), says when it is right, and
takes its lies apart one at a time.

\section{The vectorised backtest}

\begin{definition}[Backtest, vectorised backtest]\label{def:rs:vectorised-backtests:backtest}
A \emph{backtest}\index{backtest} is the simulation of a trading rule on historical data, producing the positions, trades, costs and
returns the rule would have had. A \emph{vectorised backtest}\index{vectorised backtest} computes them as whole-array operations on panels
(periods by names): target weights, returns, and costs as functions of the weight traded, without simulating orders.
\end{definition}

\begin{definition}[Decision time, execution lag, rebalance frequency]\label{def:rs:vectorised-backtests:timing}
The \emph{decision time}\index{decision time} of a book is the moment at which the information used to build it was complete. The
\emph{execution lag}\index{execution lag} is the number of periods between the decision time and the first return the book earns. The
\emph{rebalance frequency}\index{rebalance frequency} is how often the target book is recomputed and traded to.
\end{definition}

In \texttt{firm.vecbt} a book decided at the close of day $t$ with lag $L$ is held over day $t + L$: lag 0 multiplies it by the return that
produced it, lag 1 assumes it was traded at the very close it was computed from, lag 2 trades it at the next close. The book drifts with
returns between rebalances, so the weight traded at a rebalance is the target minus the drifted book, not minus the previous target;
names outside the \omterm{def:rs:universe-symbology-and-corporate-actions:universe}{tradable universe} at the \omterm{def:rs:vectorised-backtests:timing}{decision time} are not held; a name that stops listing earns its \omterm{def:rs:universe-symbology-and-corporate-actions:delisting}{delisting return} and nothing
after. The result is a \texttt{BacktestResult} (weights, trades, gross returns, costs by type, net returns, capital), the type every
backtester of the firm returns and that the performance module of chapter~22 reads.

\begin{definition}[Portfolio turnover, linear cost model]\label{def:rs:vectorised-backtests:turnover}
The \emph{portfolio turnover}\index{portfolio turnover} of a period is half the sum of the absolute weights traded, the fraction of the book
replaced. A \emph{linear cost model}\index{linear cost model} charges each trade a fixed cost per unit of weight traded (a half-spread plus
commissions), independent of its size.
\end{definition}

\section{Four fidelity levels}

\begin{definition}[Fidelity level]\label{def:rs:vectorised-backtests:fidelity}
The \emph{fidelity level}\index{fidelity level} of a \omterm{def:rs:vectorised-backtests:backtest}{backtest} is how much of the trading process it simulates. This book uses four: level~1,
the \omterm{def:rs:vectorised-backtests:backtest}{vectorised backtest} on \omterm{def:rs:market-data-for-research:bar}{bar} returns with costs as functions of turnover; level~2, the event-driven \omterm{def:rs:vectorised-backtests:backtest}{backtest} on \omterm{def:rs:market-data-for-research:bar}{bars}, with orders, fills and
a simulated clock (chapter~17); level~3, the replay of order-book messages with queue positions and latency (chapter~18); level~4, live trading
at small size or paper trading on live data (chapters 19 and 21).
\end{definition}

\begin{omfigure}
\begin{tikzpicture}[font=\small,node distance=3mm,
  lvl/.style={draw=omDef,fill=omDef!8,rounded corners=2pt,align=left,text width=3.2cm,inner sep=4pt,minimum height=2.6cm,anchor=north}]
\node[lvl] (a) at (0,0) {\textbf{Level 1}\\ vectorised\\[2pt] weights $\times$ returns;\\ costs from turnover;\\ seconds per run};
\node[lvl] (b) at (3.7,0) {\textbf{Level 2}\\ event-driven, bars\\[2pt] orders, fill models,\\ clock, calendars;\\ minutes per run};
\node[lvl] (c) at (7.4,0) {\textbf{Level 3}\\ order-book replay\\[2pt] queues, latency,\\ own impact;\\ hours per run};
\node[lvl] (d) at (11.1,0) {\textbf{Level 4}\\ live or paper\\[2pt] real fills,\\ real latency;\\ weeks per answer};
\draw[-Stealth,thick,omThm] (-1.5,-2.95) -- node[below,font=\footnotesize]{fidelity and cost of each run rise; the number of ideas that can be tested falls} (12.6,-2.95);
\end{tikzpicture}

\omcaption{The four \omterm{def:rs:vectorised-backtests:fidelity}{fidelity levels} of the firm's \omterm{def:rs:vectorised-backtests:backtest}{backtests}. A research idea moves right only if it survives on the left.}\label{fig:rs:vectorised-backtests:levels}
\end{omfigure}

The levels are not rivals: a thousand ideas are screened at level~1 in the time one is run at level~3, and a level-1 result that is negative
after costs does not need level~3 to be rejected. The question for each strategy is the lowest level whose answer the higher levels would not
change.

\section{Where the vectorised backtest is right}

It is right when trading does not change the answer: slow signals (holding periods of days to months), liquid names, books small relative
to the names' volume, and costs well approximated by a spread and a commission per unit traded. Monthly 12-1 momentum on the 500 most liquid
synthetic names trades 1.9\% of its book a day; its Sharpe ratio moves from 0.32 before costs to 0.21 after 10 basis points per unit traded,
0.19 after borrow and 0.17 when traded at the next close (\cref{fig:rs:vectorised-backtests:waterfall}). None of those steps depends on how the
trades are executed, and a level-2 or level-3 \omterm{def:rs:vectorised-backtests:backtest}{backtest} would add little. It is wrong, in the sense of missing the answer, when fills depend on the
order book (passive orders, chapter~17), when the strategy's own trades move prices (chapter~27), when timing within the day matters, and when
events such as halts, auctions or corporate actions intervene between the decision and the trade.

\section{The lies it tells}

\begin{proposition}[Costs against turnover]\label{prop:rs:vectorised-backtests:breakeven}
A book with mean gross return $g$ per period and turnover $\tau$ per period (half the weight traded) breaks even at a linear cost of
$g/(2\tau)$ per unit of weight traded.
\end{proposition}

\begin{proof}
The weight traded per period is $2\tau$, and at a cost $c$ per unit the cost per period is $2\tau c$; set it equal to $g$.
\end{proof}

The one-day reversal on the synthetic market, rebuilt daily as a dollar-neutral book of gross exposure 1, corrected one lie at a time
(\texttt{rs\_vecbt.waterfall}; each step keeps the corrections before it):

\begin{center}
\small
\begin{tabular}{@{}lccc@{}}
\toprule
step & Sharpe ratio & mean return a year & turnover a day\\
\midrule
naive (survivors, trade at the signal's close) & 7.17 & 27.0\% & 74\%\\
names listed at each decision & 7.22 & 27.4\% & 74\%\\
the 500 most liquid names & 6.95 & 27.4\% & 75\%\\
10 bp per unit traded & $-2.58$ & $-10.2\%$ & 75\%\\
borrow at 50 bp a year & $-2.64$ & $-10.4\%$ & 75\%\\
trade at the next close & $-9.45$ & $-37.7\%$ & 75\%\\
\bottomrule
\end{tabular}
\end{center}

\textbf{The universe.} The naive universe is the list of names alive at the end of the sample, the one a researcher downloads today. For the
reversal it changes little (7.17 against 7.22 on the point-in-time list), for momentum a great deal, and in the direction a textbook does not
predict: 0.21 against 0.44 (\cref{fig:rs:vectorised-backtests:waterfall}). The names that later delisted were momentum's losers, and the
short side profited from their final falls; \omterm{def:rs:point-in-time-data-and-the-biases:survivorship}{survivorship bias} removes that profit. Its sign depends on the strategy; its size can only be measured.

\textbf{Costs.} The reversal's gross return is large and its turnover larger: 74\% of the book a day. At 10 basis points per unit of weight traded
the cost is $2 \times 0.75 \times 0.001 \times 252 = 38\%$ a year against a gross 27\%, and the Sharpe ratio goes from 6.95 to $-2.58$. By
\cref{prop:rs:vectorised-backtests:breakeven} the book breaks even at 7.3 basis points. Novy-Marx and Velikov (2016) found the same on real
anomalies: most strategies with less than 50\% monthly turnover still earned significant spreads after costs when designed to limit them, and few
with higher turnover did.

\textbf{Borrow.} Short positions pay a fee to the lender of the stock (Book~1, chapter~6), 50 basis points a year on general collateral and far more
on hard-to-borrow names; here it costs 0.25\% a year and little Sharpe ratio. It matters more for strategies whose shorts concentrate in small,
crowded names, which the reversal's do not.

\textbf{Timing.} A book computed from a closing price cannot be traded at that price: the closing auction's orders are due before the close is known
(Book~1, chapter~13). Traded at the next close instead, the reversal has no alpha left (its information lasts a day, chapter~13) and keeps its
costs: $-9.45$. The honest level-1 answer would use a signal computed from prices known before the close and a cost model of the closing auction;
the \omterm{def:rs:vectorised-backtests:backtest}{vectorised backtest} can hold that assumption, but only if someone writes it down.

\textbf{Compounding.} Summing daily returns instead of compounding them is harmless for small returns and not for large ones. The costed reversal's
capital, compounded, falls to 0.35 of its start over ten years; summed, it reaches $-0.04$, a capital no book can have.

\textbf{Look-ahead.} Multiplying today's book by today's return makes the reversal lose 694\% a year ($-65$ Sharpe ratio): its signal is minus that
return. Look-ahead does not always flatter; it always lies.

\begin{omfigure}
\begin{tikzpicture}
\begin{axis}[name=rev,width=6.8cm,height=5.6cm,ybar,bar width=10pt,title={\small one-day reversal},ylabel={Sharpe ratio},tick label style={font=\footnotesize},
  label style={font=\small},xtick={0,1,2,3,4,5},xticklabels={naive,PIT,liquid,costs,borrow,next},x tick label style={rotate=45,anchor=east},
  ymin=-10.5,ymax=8.5,grid=major,grid style={gray!25},table/col sep=comma]
\addplot[fill=omThm!45,draw=omThm] table[x=k,y=reversal]{figdata/research/16-vectorised-backtests/waterfall.csv};
\end{axis}
\begin{axis}[at={(rev.east)},anchor=west,xshift=1.6cm,width=6.8cm,height=5.6cm,ybar,bar width=10pt,title={\small momentum 12-1, monthly},
  tick label style={font=\footnotesize},label style={font=\small},xtick={0,1,2,3,4,5},xticklabels={naive,PIT,liquid,costs,borrow,next},
  x tick label style={rotate=45,anchor=east},ymin=0,ymax=0.5,grid=major,grid style={gray!25},table/col sep=comma]
\addplot[fill=omDef!55,draw=omDef] table[x=k,y=momentum]{figdata/research/16-vectorised-backtests/waterfall.csv};
\end{axis}
\end{tikzpicture}

\omcaption{Sharpe ratio of two strategies on \texttt{firm.synthmkt} as the \omterm{def:rs:vectorised-backtests:backtest}{backtest}'s lies are corrected one at a time: survivors' universe (naive),
point-in-time universe (PIT), 500 most liquid names, 10 bp per unit traded, borrow at 50 bp a year, trading at the next close. Note the two
scales.}\label{fig:rs:vectorised-backtests:waterfall}
\end{omfigure}

\begin{omfigure}
\begin{tikzpicture}
\begin{axis}[width=11cm,height=5.6cm,xlabel={years},ylabel={capital ($\log_{10}$)},tick label style={font=\small},label style={font=\small},
  xmin=0,xmax=10,ymin=-2,ymax=1.5,grid=major,grid style={gray!25},legend columns=2,
  legend style={font=\footnotesize,at={(0.5,-0.25)},anchor=north,draw=none,fill=none,
  /tikz/every even column/.append style={column sep=8pt}},table/col sep=comma]
\addplot[black!50,thick,dashed] table[x=year,y=gross]{figdata/research/16-vectorised-backtests/capital.csv};
\addplot[omThm,thick] table[x=year,y=costed]{figdata/research/16-vectorised-backtests/capital.csv};
\addplot[omThm,thick,dotted] table[x=year,y=next]{figdata/research/16-vectorised-backtests/capital.csv};
\addplot[omDef,very thick] table[x=year,y=momentum]{figdata/research/16-vectorised-backtests/capital.csv};
\legend{reversal before costs,reversal after costs and borrow,reversal traded at the next close,momentum after all corrections}
\end{axis}
\end{tikzpicture}

\omcaption{Compounded capital of the reversal on the 500 most liquid synthetic names before costs, after costs and borrow, and traded at the next close,
with monthly momentum after all corrections. Data: \texttt{rs\_vecbt} on \texttt{firm.synthmkt}.}\label{fig:rs:vectorised-backtests:capital}
\end{omfigure}

\section{A checklist}

Every level-1 \omterm{def:rs:vectorised-backtests:backtest}{backtest} the firm runs states, and \texttt{firm.vecbt} enforces where it can:
\begin{enumerate}
\item the \omterm{def:rs:vectorised-backtests:timing}{decision time} of every input, and the \omterm{def:rs:vectorised-backtests:timing}{execution lag} from it to the trade (never 0; 1 only with a stated reason);
\item the \omterm{def:rs:universe-symbology-and-corporate-actions:universe}{tradable universe} at each decision, point-in-time, with \omterm{def:rs:universe-symbology-and-corporate-actions:delisting}{delisting returns};
\item the cost per unit traded, by name if liquidity differs, and the break-even cost (\cref{prop:rs:vectorised-backtests:breakeven});
\item borrow fees and availability for the shorts, and the financing of any leverage and the interest on cash;
\item the \omterm{def:rs:vectorised-backtests:timing}{rebalance frequency} and the drift of the book between rebalances;
\item compounding, and capital that cannot go below zero;
\item the gross exposure, and position caps against liquidity;
\item the number of variants tried before this one (chapter~20), and the level at which the result must be confirmed.
\end{enumerate}
Arnott, Harvey and Markowitz (2019) set out such a protocol for research with machine learning, stressing that its methods need far more data than
finance has.

\section{Tutorial: Sharpe 7.2 on Monday}\label{tut:rs:vectorised-backtests:tut}

\begin{tutorial}
\textbf{Goal.} \omterm{def:rs:vectorised-backtests:backtest}{Backtest} the one-day reversal and monthly momentum on \texttt{firm.synthmkt} in one vectorised pass each, correct the lies one at a
time, and compute the reversal's break-even cost. \textbf{End state:}
\cref{fig:rs:vectorised-backtests:waterfall,fig:rs:vectorised-backtests:capital}; the table.
\begin{tutsteps}
\item \textbf{The core loop}: the lagged book, the trade from the drifted book, costs, borrow, financing.
\omcode{code/firm/vecbt/firm_vecbt.py}{90}{101}{The level-1 backtest's period loop.}\label{lst:rs:vectorised-backtests:loop}
\item \textbf{The \omterm{def:rs:universe-symbology-and-corporate-actions:universe}{tradable universe}}: the most liquid names by trailing dollar volume, as known at each close.
\omcode{code/research/16-vectorised-backtests/python/rs_vecbt.py}{37}{43}{A point-in-time liquid universe.}\label{lst:rs:vectorised-backtests:liquid}
\item \textbf{Run} \texttt{waterfall()} for both signals, \texttt{lookahead()}, \texttt{breakeven\_cost()}, \texttt{compounding()} and
\texttt{fig\_vecbt.py}.
\end{tutsteps}
\textbf{What to change next.} Smooth the reversal signal over three days and watch turnover and the break-even cost; charge costs by name in proportion to
each name's volatility.
\end{tutorial}

\section{Build: the level-1 backtester}\label{bld:rs:vectorised-backtests:vecbt}

\begin{build}
\bfield{Purpose} The firm's screening backtester: seconds per run, honest by construction on timing, universe, costs and financing, and the source of the
\texttt{BacktestResult} type every later backtester returns.
\bfield{Interface} \texttt{BacktestResult(dates, names, weights, trades, gross, costs, net, capital, meta)} with \texttt{turnover},
\texttt{gross\_exposure}, \texttt{net\_exposure}; \texttt{signal\_to\_weights(signal, universe, gross, neutral)}; \texttt{backtest(weights, returns, lag,
universe, cost, borrow, cash\_rate, borrow\_spread, cap, compound, periods)}.
\bfield{Rules} The lag is an argument, not an assumption; the universe is applied at the \omterm{def:rs:vectorised-backtests:timing}{decision time}; trades are measured from the drifted book; costs
are charged on every trade; capital compounds by default.
\bfield{Acceptance tests} \texttt{code/firm/vecbt/tests/}: lags 0, 1 and 2 on a hand series; compounding against summing; the trade from a drifted book;
costs and turnover of an alternating book; borrow, cash interest and leverage financing; universe masks and caps; dollar-neutral weights of a given gross.
\bfield{Stretch} Costs from a square-root impact model (chapter~27); borrow availability by name; multi-currency cash.
\end{build}

\begin{omsources}
\item R.~Novy-Marx and M.~Velikov, ``A taxonomy of anomalies and their trading costs'', \emph{Review of Financial Studies} 29(1), 2016.
\item R.~Arnott, C.~R.~Harvey and H.~Markowitz, ``A backtesting protocol in the era of machine learning'', \emph{Journal of Financial Data Science} 1(1), 2019.
\end{omsources}

\section{Exercises}

\begin{exercise}[$\star$]\label{exo:rs:vectorised-backtests:1}
A book earns 12\% a year gross and turns over 20\% of itself a day (252 days). At what cost per unit traded does it break even?
\end{exercise}

\begin{exercise}[$\star$]\label{exo:rs:vectorised-backtests:2}
A book holds 60\% in A and 40\% in B. A returns 10\% and B 0\% over the period. What are the drifted weights, and what must be traded to return to
60/40?
\end{exercise}

\begin{exercise}[$\star$]\label{exo:rs:vectorised-backtests:3}
Daily returns of $+50\%$ and $-40\%$ alternate for ten days. What is the capital after ten days compounded, and what does summing the returns suggest?
\end{exercise}

\begin{exercise}[$\star\star$]\label{exo:rs:vectorised-backtests:4}
Why does the look-ahead \omterm{def:rs:vectorised-backtests:backtest}{backtest} of the reversal lose money while that of a momentum signal including today's return would make it?
\end{exercise}

\begin{exercise}[$\star\star$]\label{exo:rs:vectorised-backtests:5}
Why can \omterm{def:rs:point-in-time-data-and-the-biases:survivorship}{survivorship bias} raise one strategy's \omterm{def:rs:vectorised-backtests:backtest}{backtest} and lower another's? Give the mechanism for momentum in the synthetic market.
\end{exercise}

\begin{exercise}[$\star\star$]\label{exo:rs:vectorised-backtests:6}
A dollar-neutral book of gross exposure 2 posts its capital as collateral earning 3\% a year and pays 50 basis points a year to borrow its shorts. What
does financing add or cost per year, per unit of capital?
\end{exercise}

\begin{exercise}[$\star\star\star$]\label{exo:rs:vectorised-backtests:7}
\emph{Coding.} Rebuild the reversal from the average of the last three days' returns instead of the last one. How do turnover, the gross Sharpe ratio and
the break-even cost change?
\end{exercise}

\begin{exercise}[$\star\star\star$]\label{exo:rs:vectorised-backtests:8}
\emph{Find the flaw.} ``Our strategy trades at the close using the closing price, which is fine because we use market-on-close orders.''
\end{exercise}

\section{Problem: Sharpe 7.2 on Monday, $-2.6$ on Friday}

\begin{problem}[{Weekend problem --- a backtest, lie by lie}]\label{pb:rs:vectorised-backtests:1}
The one-day reversal and monthly 12-1 momentum on \texttt{firm.synthmkt}, dollar neutral, gross exposure 1.

\medskip\noindent\textbf{Part I --- The naive \omterm{def:rs:vectorised-backtests:backtest}{backtest}.}
\begin{enumerate}
\item What Sharpe ratio and mean return does the naive reversal \omterm{def:rs:vectorised-backtests:backtest}{backtest} show, and what does it assume?
\item What does the lag-0 version show, and why is its sign negative?
\item What turnover does the reversal have?
\item Which inputs of the naive \omterm{def:rs:vectorised-backtests:backtest}{backtest} could a researcher have known at the \omterm{def:rs:vectorised-backtests:timing}{decision time}?
\end{enumerate}

\medskip\noindent\textbf{Part II --- The universe.}
\begin{enumerate}[resume]
\item What do the point-in-time and liquid universes do to each strategy?
\item Why does \omterm{def:rs:point-in-time-data-and-the-biases:survivorship}{survivorship bias} lower momentum here?
\item What would it do to a strategy that buys cheap, distressed stocks?
\end{enumerate}

\medskip\noindent\textbf{Part III --- Costs and timing.}
\begin{enumerate}[resume]
\item What do 10 basis points per unit traded do to the reversal and to momentum?
\item What is the reversal's break-even cost?
\item What does borrow cost?
\item What happens when the reversal is traded at the next close, and why?
\item What does compounding change for the costed reversal?
\end{enumerate}

\medskip\noindent\textbf{Part IV --- The verdict.}
\begin{enumerate}[resume]
\item State the \emph{named result}: the Sharpe waterfall from the naive to the honest level-1 \omterm{def:rs:vectorised-backtests:backtest}{backtest}, lie by lie, for both strategies.
\item Which strategy is a level-1 result, and which needs a higher \omterm{def:rs:vectorised-backtests:fidelity}{fidelity level} before anyone believes it?
\item What would a closing-auction cost model and a signal from 15:50 prices change?
\item How would you cut the reversal's turnover without losing its alpha?
\item Which line of the checklist did the naive \omterm{def:rs:vectorised-backtests:backtest}{backtest} violate first?
\item What would Novy-Marx and Velikov's evidence lead you to expect for a real daily reversal?
\item What should the \omterm{def:rs:the-research-process:log}{research log} record about this \omterm{def:rs:vectorised-backtests:backtest}{backtest}?
\item In one sentence: what is a \omterm{def:rs:vectorised-backtests:backtest}{vectorised backtest} good for?
\end{enumerate}
\end{problem}

\section{Interview questions}

\begin{interviewq}[$\star$ \iqroles{researcher}]\label{iq:rs:vectorised-backtests:1}
Name five ways a \omterm{def:rs:vectorised-backtests:backtest}{backtest} can overstate performance.
\end{interviewq}

\begin{interviewq}[$\star\star$ \iqroles{researcher, developer}]\label{iq:rs:vectorised-backtests:2}
Write, in words, the steps of a \omterm{def:rs:vectorised-backtests:backtest}{vectorised backtest} of a daily long-short strategy, and where each bias enters.
\end{interviewq}

\begin{interviewq}[$\star\star$ \iqroles{researcher, trader}]\label{iq:rs:vectorised-backtests:3}
A strategy's gross Sharpe ratio is 7. Why might you still not trade it?
\end{interviewq}

\begin{interviewq}[$\star\star$ \iqroles{researcher}]\label{iq:rs:vectorised-backtests:4}
When is a \omterm{def:rs:vectorised-backtests:backtest}{vectorised backtest} good enough, and when do you need an event-driven one?
\end{interviewq}

\begin{interviewq}[$\star\star$ \iqroles{researcher, risk}]\label{iq:rs:vectorised-backtests:5}
How do you account for borrow and financing in a long-short \omterm{def:rs:vectorised-backtests:backtest}{backtest}?
\end{interviewq}

\begin{interviewq}[$\star\star\star$ \iqroles{researcher}]\label{iq:rs:vectorised-backtests:6}
Derive the break-even cost of a strategy from its gross return and turnover, and use it to compare a daily reversal with monthly momentum.
\end{interviewq}
